% Capítulo 3: Proposta e Metodologia
\chapter{Proposta e metodologia}\label{chp:desenvolvimento}

% \section{Framework arquitetural proposto}\label{sec:framework_arquitetural}

% Com base na análise sistemática do estado da arte realizada no
% \Capitulo{chp:levantamento}, propõe-se um framework arquitetural heterogêneo que
% integra as melhores técnicas identificadas em um sistema único e otimizado.

% \subsection{Justificativa e requisitos}\label{subsec:justificativa_requisitos}

% A justificativa para o framework proposto baseia-se em três pilares fundamentais
% identificados na revisão sistemática:

% \textbf{1. Lacuna de Integração}: Apenas 1 trabalho (\textcite{Shin2020}) demonstra sistema completo integrado, enquanto 96\% dos trabalhos focam em técnicas isoladas. Esta lacuna limita significativamente a aplicabilidade prática das soluções desenvolvidas.

% \textbf{2. Potencial Combinado}: A análise quantitativa revela que técnicas individuais já atingem performance adequada:
% \begin{itemize}
%     \item \textcite{Martins2019}: 10,000 fps com 99.7\% precisão (FPGA)
%     \item \textcite{Diaz2019}: 330 fps com compressão >20:1 (GPU)
%     \item \textcite{Hwang2011}: 43.5x melhoria energética (codesign)
% \end{itemize}

% \textbf{3. Necessidade de Trade-offs Holísticos}: A ausência de análise integrada de trade-offs entre precisão, latência e consumo energético impede a otimização global do sistema.

% Os requisitos do framework incluem:
% \begin{itemize}
%     \item \textbf{Performance}: Throughput $\geq$100 fps para aplicações em tempo real
%     \item \textbf{Latência}: <50ms end-to-end por frame
%     \item \textbf{Consumo}: <15W total para aplicações embarcadas
%     \item \textbf{Precisão}: >95\% mantendo robustez
%     \item \textbf{Escalabilidade}: Suporte a diferentes resoluções e bandas
% \end{itemize}

% \subsection{Especificações de hardware e software}
% \label{subsec:especificacoes_hw_sw}

% \textbf{Arquitetura Hardware Proposta}:
% \begin{itemize}
%     \item \textbf{FPGA}: Xilinx Zynq UltraScale+ (pré-processamento especializado)
%     \item \textbf{GPU}: NVIDIA Jetson Orin Nano (processamento principal)
%     \item \textbf{CPU}: ARM Cortex-A78 (classificação e controle)
%     \item \textbf{Memória}: 16GB LPDDR5 compartilhada
%     \item \textbf{Comunicação}: PCIe Gen4 + shared memory
% \end{itemize}

% \textbf{Justificativa da Escolha}:
% \begin{itemize}
%     \item \textbf{FPGA}: Baseada em \textcite{Hwang2011} para baixa latência e alta eficiência
%     \item \textbf{GPU}: Baseada em \textcite{Diaz2019} para processamento paralelo massivo
%     \item \textbf{CPU}: Baseada em \textcite{TakuNet2024} para classificação leve e controle adaptativo
% \end{itemize}

% \section{Pipeline de processamento}\label{sec:pipeline_processamento}

% O pipeline proposto segue arquitetura de três estágios, cada um otimizado para
% funções específicas baseadas nas melhores práticas identificadas na literatura.

% \subsection{Pré-processamento no FPGA}\label{subsec:preprocessamento_fpga}

% \textbf{Inspiração}: \textcite{Hwang2011} identificou que o pré-processamento consome 91\% do tempo computacional total.

% \textbf{Técnicas Implementadas}:

% \textbf{1. Correção Radiométrica ELM}: Baseada em \textcite{Shin2024}, implementa Empirical Line Method para correção de reflectância, resultando em melhoria de 5-55\% na precisão radiométrica.

% \textbf{2. Seleção de Bandas EMCR}: Baseada em \textcite{Martins2019}, implementa seleção inteligente de bandas resultando em redução de 80\% dos dados mantendo 99.7\% de precisão.

% \textbf{3. Compressive Sensing Encoder}: Baseada em \textcite{Lim2022}, implementa encoder CS para redução adicional de 50-70\% dos dados.

% \textbf{Recursos FPGA Estimados}:
% \begin{itemize}
%     \item DSPs: 45\% (600 de 1340 disponíveis)
%     \item BRAMs: 60\% (850 de 1418 disponíveis)
%     \item LUTs: 70\% (240K de 343K disponíveis)
% \end{itemize}

% \subsection{Núcleo de processamento na GPU}
% \label{subsec:nucleo_processamento_gpu}

% \textbf{Inspiração}: \textcite{Diaz2019} demonstrou 330 fps em Jetson TX2 com algoritmo HyperLCA.

% \textbf{Técnicas Implementadas}:

% \textbf{1. Decompressor CS CUDA}: Baseada em \textcite{Lim2022}, implementa algoritmo CGNE otimizado para reconstrução de dados comprimidos com complexidade O(n³) mas altamente paralelizável.

% \textbf{2. Processamento Multiescala CNN}: Baseada em \textcite{TakuNet2024}, implementa arquitetura depth-wise com apenas 37.685 parâmetros, atingindo >650 fps com F1-score de 94.3\%.

% \textbf{3. Pipeline de Fusão Espacial-Espectral}: Baseada em \textcite{UAV2024}, implementa mecanismos de atenção para fusão otimizada de características espaciais e espectrais.

% \textbf{Otimizações GPU}:
% \begin{itemize}
%     \item FP16 precision (redução 50\% memória)
%     \item Tensor cores utilization
%     \item Memory coalescing optimization
%     \item Warp-level primitives
% \end{itemize}

% \subsection{Classificação e controle na CPU}
% \label{subsec:classificacao_controle_cpu}

% \textbf{Inspiração}: \textcite{Martins2019} demonstrou 0.1ms/pixel com SVM otimizado em FPGA.

% \textbf{Técnicas Implementadas}:

% \textbf{1. Classificador SVM Otimizado}: Implementa classificação com distância de Hamming para máxima eficiência, baseado na arquitetura de \textcite{Martins2019}.

% \textbf{2. Controle Adaptativo de Qualidade}: Sistema inteligente que ajusta dinamicamente parâmetros de compressão e processamento baseado na carga computacional atual.

% \textbf{3. Gestão Inteligente de Energia}: Baseada em \textcite{Hwang2011}, implementa power gating e ajuste dinâmico de frequências.

% \subsection{Interfaces, buffers e sincronização}
% \label{subsec:interfaces_buffers}

% \textbf{Arquitetura de Comunicação}:
% \begin{itemize}
%     \item \textbf{FPGA $\leftrightarrow$ GPU}: PCIe Gen4 (16 GB/s) com DMA otimizado
%     \item \textbf{GPU $\leftrightarrow$ CPU}: Shared memory (unified memory) para baixa latência
%     \item \textbf{CPU $\leftrightarrow$ Sistema}: Ethernet/WiFi para decisões e controle
% \end{itemize}

% \textbf{Sincronização}:
% \begin{itemize}
%     \item \textbf{DMA transfers} entre FPGA e memória compartilhada
%     \item \textbf{CUDA streams} para overlap computation/communication
%     \item \textbf{Lock-free queues} para comunicação CPU-GPU
% \end{itemize}

% \section{Plano de avaliação experimental}\label{sec:plano_avaliacao}

% \subsection{Datasets e cenários de teste}\label{subsec:datasets_cenarios}

% \textbf{Datasets de Referência} (baseados em \textcite{Lou2024} e \textcite{Ullah2020}):
% \begin{itemize}
%     \item \textbf{AVIRIS Indian Pines}: Dataset agrícola de referência (145×145×200 bandas)
%     \item \textbf{Pavia University}: Dataset urbano de alta resolução (610×340×103 bandas)
%     \item \textbf{Salinas Valley}: Dataset agrícola diversificado (512×217×204 bandas)
%     \item \textbf{Dataset próprio UAV}: Validação em aplicação agrícola real
% \end{itemize}

% \textbf{Cenários de Teste}:
% \begin{enumerate}
%     \item \textbf{Agricultura de Precisão}: UAV a 100m altitude, área 1 hectare, detecção de estresse hídrico
%     \item \textbf{Monitoramento Urbano}: Classificação LULC em tempo real, processamento contínuo 30 fps
%     \item \textbf{Aplicação de Emergência}: Detecção de incêndios florestais, latência crítica <100ms
% \end{enumerate}

% \subsection{Métricas e protocolos de medição}\label{subsec:metricas_protocolos}

% \textbf{Métricas de Performance}:
% \begin{itemize}
%     \item \textbf{Throughput}: FPS médio e pico, variação temporal
%     \item \textbf{Latência}: End-to-end, por estágio, percentil 95\%
%     \item \textbf{Throughput}: MB/s, frames/segundo, pixels/segundo
% \end{itemize}

% \textbf{Métricas de Eficiência}:
% \begin{itemize}
%     \item \textbf{Energia}: Joules por frame, watts médios e pico
%     \item \textbf{Eficiência}: fps/W, MB/J, pixels/J
%     \item \textbf{Distribuição}: Consumo por módulo, overhead de comunicação
% \end{itemize}

% \textbf{Métricas de Qualidade}:
% \begin{itemize}
%     \item \textbf{Precisão}: Accuracy, F1-score, RMSE
%     \item \textbf{Robustez}: Tolerância a falhas, degradação graciosa
%     \item \textbf{Consistência}: Variação temporal, dependência de condições
% \end{itemize}

% \subsection{Benchmarks de referência}\label{subsec:benchmarks_referencia}

% \textbf{Baseline Single-Core}: Implementação CPU pura para comparação de performance e eficiência.

% \textbf{Implementações Isoladas}: Comparação com trabalhos individuais da literatura:
% \begin{itemize}
%     \item \textcite{Martins2019}: FPGA SVM (10,000 fps, 0.1ms)
%     \item \textcite{Diaz2019}: GPU HyperLCA (330 fps, 3.03ms)
%     \item \textcite{TakuNet2024}: CPU CNN (650 fps, 1.54ms)
% \end{itemize}

% \textbf{Sistemas Completos}: Comparação com \textcite{Shin2020} (UAV integrado, 25 fps, 40ms).

% \section{Metas quantitativas e critérios de sucesso}\label{sec:metas_criterios}

% \subsection{Targets de Performance}\label{subsec:targets_performance}

% \textbf{Throughput}: $\geq$100 fps para resolução 614$\times$512$\times$224 bandas
% \begin{itemize}
%     \item \textbf{Justificativa}: Baseado em \textcite{Diaz2019} (330 fps) e \textcite{TakuNet2024} (650 fps)
%     \item \textbf{Target conservador}: Considerando overhead de integração
% \end{itemize}

% \textbf{Latência}: <50ms end-to-end por frame
% \begin{itemize}
%     \item \textbf{Justificativa}: Baseado em \textcite{Lim2022} (<10ms CS) e \textcite{Martins2019} (0.1ms classificação)
%     \item \textbf{Overhead}: Comunicação entre módulos e sincronização
% \end{itemize}

% \textbf{Consumo}: <15W total
% \begin{itemize}
%     \item \textbf{Justificativa}: Baseado em \textcite{Hwang2011} (43.5x melhoria energética)
%     \item \textbf{Plataforma embarcada}: Restrições de potência para aplicações UAV
% \end{itemize}

% \subsection{Critérios de Sucesso}\label{subsec:criterios_sucesso}

% \textbf{Critérios Técnicos}:
% \begin{itemize}
%     \item \textbf{Performance}: Throughput $\geq$100 fps, latência <50ms
%     \item \textbf{Eficiência}: Consumo <15W, eficiência >6.7 fps/W
%     \item \textbf{Qualidade}: Precisão >95\%, robustez >90\%
% \end{itemize}

% \textbf{Critérios de Integração}:
% \begin{itemize}
%     \item \textbf{Modularidade}: Estágios independentes e intercambiáveis
%     \item \textbf{Escalabilidade}: Suporte a diferentes resoluções
%     \item \textbf{Configurabilidade}: Trade-offs dinâmicos precisão vs recursos
% \end{itemize}

% \section{Riscos técnicos e estratégias de mitigação}
% \label{sec:riscos_estrategias}

% \subsection{Riscos Identificados}\label{subsec:riscos_identificados}

% \textbf{1. Complexidade de Integração}:
% \begin{itemize}
%     \item \textbf{Risco}: Overhead de comunicação entre módulos pode comprometer ganhos
%     \item \textbf{Mitigação}: Otimização de interfaces, uso de shared memory, DMA otimizado
% \end{itemize}

% \textbf{2. Sincronização de Pipeline}:
% \begin{itemize}
%     \item \textbf{Risco}: Gargalos em um estágio podem afetar todo o sistema
%     \item \textbf{Mitigação}: Buffers adaptativos, balanceamento dinâmico de carga
% \end{itemize}

% \textbf{3. Consumo Energético Integrado}:
% \begin{itemize}
%     \item \textbf{Risco}: Soma dos consumos pode exceder limite de 15W
%     \item \textbf{Mitigação}: Power gating, ajuste dinâmico de frequências, otimizações por estágio
% \end{itemize}

% \subsection{Estratégias de Mitigação}\label{subsec:estrategias_mitigacao}

% \textbf{1. Prototipagem Incremental}:
% \begin{itemize}
%     \item Implementação e validação de cada módulo isoladamente
%     \item Integração gradual com validação contínua
%     \item Identificação precoce de problemas de integração
% \end{itemize}

% \textbf{2. Simulação e Modelagem}:
% \begin{itemize}
%     \item Modelagem matemática do pipeline antes da implementação
%     \item Simulação de diferentes cenários de carga
%     \item Análise de trade-offs teóricos
% \end{itemize}

% \textbf{3. Validação Contínua}:
% \begin{itemize}
%     \item Testes automatizados em cada estágio de desenvolvimento
%     \item Métricas de qualidade em tempo real
%     \item Feedback loop para otimizações
% \end{itemize}

% \section{Cronograma de desenvolvimento}\label{sec:cronograma_desenvolvimento}

% \subsection{Fase 1: Análise e Baseline (Concluída)}\label{subsec:fase1_analise}

% \textbf{Período}: Agosto-Dezembro 2025
% \begin{itemize}
%     \item \textbf{Concluído}: Revisão sistemática de 25 artigos
%     \item \textbf{Concluído}: Catalogação quantitativa de técnicas
%     \item \textbf{Concluído}: Identificação de gaps e oportunidades
%     \item \textbf{Concluído}: Framework arquitetural conceitual
% \end{itemize}

% \subsection{Fase 2: Desenvolvimento Experimental (Pós-qualificação)}
% \label{subsec:fase2_desenvolvimento}

% \textbf{Período}: Janeiro-Junho 2026 (6 meses)

% \textbf{Meses 1-2: Análise e Baseline}
% \begin{itemize}
%     \item Setup ambiente heterogêneo (FPGA + GPU + CPU)
%     \item Profiling sistemático de algoritmos base
%     \item Implementação baseline single-core
%     \item Medições detalhadas e relatório
% \end{itemize}

% \textbf{Meses 3-6: Implementação Core}
% \begin{itemize}
%     \item \textbf{Mês 3}: Módulo FPGA preprocessing
%     \item \textbf{Mês 4}: Módulo GPU processing
%     \item \textbf{Mês 5}: Módulo CPU classification
%     \item \textbf{Mês 6}: Integração inicial dos módulos
% \end{itemize}

% \subsection{Fase 3: Integração e Validação (Pós-qualificação)}
% \label{subsec:fase3_integracao}

% \textbf{Período}: Julho-Setembro 2026 (3 meses)

% \textbf{Meses 7-9: Integração e Otimização}
% \begin{itemize}
%     \item \textbf{Mês 7}: Pipeline heterogêneo completo
%     \item \textbf{Mês 8}: Otimizações de comunicação
%     \item \textbf{Mês 9}: Balanceamento e power management
% \end{itemize}

% \textbf{Meses 10-11: Validação e Refinamento}
% \begin{itemize}
%     \item \textbf{Mês 10}: Validação experimental extensiva
%     \item \textbf{Mês 11}: Auto-tuning e otimizações finais
% \end{itemize}

% \section{Considerações finais do capítulo}\label{sec:consideracoes_finais_cap3}

% O framework arquitetural heterogêneo proposto representa uma abordagem integrada
% que combina as melhores técnicas identificadas na revisão sistemática do estado
% da arte. A arquitetura CPU+GPU+FPGA oferece flexibilidade e eficiência para
% diferentes estágios de processamento.

% \textbf{Principais Inovações}:
% \begin{itemize}
%     \item \textbf{Integração Sistemática}: Primeira proposta que integra técnicas comprovadas em pipeline único
%     \item \textbf{Otimização Holística}: Análise de trade-offs considerando todo o sistema
%     \item \textbf{Validação Robusta}: Metodologia de avaliação com múltiplas métricas
% \end{itemize}

% \textbf{Contribuições Esperadas}:
% \begin{itemize}
%     \item \textbf{Teórica}: Framework conceitual para sistemas heterogêneos
%     \item \textbf{Metodológica}: Protocolos de validação padronizados
%     \item \textbf{Prática}: Especificações para implementação
% \end{itemize}

% \textbf{Próximos Passos}:
% \begin{itemize}
%     \item Preparação para qualificação UNICAMP
%     \item Implementação experimental pós-qualificação
%     \item Validação em cenários reais de aplicação
% \end{itemize}

% Este framework estabelece as bases para o desenvolvimento de sistemas
% hiperespectrais embarcados de próxima geração, capazes de atender aos requisitos
% rigorosos de aplicações em tempo real com eficiência energética superior ao
% estado da arte atual.
